% Source: https://artofproblemsolving.com/community/c3h2000015_expected_value_problem_tell_me_if_in_wrong_forum
% Stable identity: aops-topic-2000015-q01
Let $X$ be the score. Then $E[X]=\sum_{n\ge1}P(X\ge n)$, and $X\ge n$ means the first $n$ darts all hit.

\textbf{One hit.} Given the current radius $r$, a dart hits with probability $\dfrac{\pi r^2}{4}$ (area of the disc over the area $4$ of the square). Given a hit, the dart is uniform in the disc, and then $u=\dfrac{d^2}{r^2}$ is uniform on $[0,1]$ (because $P(d\le t)=\dfrac{t^2}{r^2}$). The new squared radius is $r^2-d^2=r^2(1-u)$, and $1-u$ is again uniform on $[0,1]$.

So if the first $k$ darts hit, $r_k^2=U_1U_2\cdots U_k$ with independent uniform $U_i$ on $[0,1]$, and
\[
P(X\ge n)=E\left[\prod_{k=0}^{n-1}\dfrac{\pi r_k^2}{4}\right]=\left(\dfrac{\pi}{4}\right)^n E\left[U_1^{\,n-1}U_2^{\,n-2}\cdots U_{n-1}^{\,1}\right].
\]
(The $k$-th factor is the probability that dart $k+1$ hits given the previous hits, and $r_0^2=1$.) Since $E[U^m]=\dfrac{1}{m+1}$ and the $U_i$ are independent,
\[
E\left[U_1^{\,n-1}\cdots U_{n-1}^{\,1}\right]=\dfrac{1}{n}\cdot\dfrac{1}{n-1}\cdots\dfrac{1}{2}=\dfrac{1}{n!},\qquad P(X\ge n)=\dfrac{1}{n!}\left(\dfrac{\pi}{4}\right)^n .
\]
Therefore
\[
E[X]=\sum_{n\ge1}\dfrac{1}{n!}\left(\dfrac{\pi}{4}\right)^n=\boxed{e^{\dfrac{\pi}{4}}-1}\approx1.193 .
\]
(If the final, missed throw is also counted, as in some versions of this puzzle, the expected number of throws is $e^{\dfrac{\pi}{4}}$.)
